\runninghead{The Kingdom Handed Over}
\properlane{kingdom-handed-over}{The Kingdom Handed Over That Has No End}

The end of Christ's kingship is the Father. The eternal kingdom that the
Preface says he hands over is his redeemed people, whom he brings to the
vision of God, so that handing over the kingdom and reigning for ever are one
consummation and not an abdication. The session for ever that the Communion
sings begins already in the souls he has cleansed and is completed after the
judgment. The reading begins from the Preface, which ends the Mass's
thanksgiving with the kingdom's final destination, and lets the Alleluia and
the Communion say what kind of end that is.

\subsection{A kingdom delivered up that shall not be taken away}
\label{sec:handed-question}

The Preface follows St Paul's order in 1 Cor 15. The Son is anointed priest
and king; he offers himself on the cross; all creatures are subjected to his
rule; and then he hands over \latin{aeternum et universale regnum} to the
Father's majesty. The Apostle's own words stand behind the verb: ``Afterwards
the end: when he shall have delivered up the kingdom to God and the Father:
when he shall have brought to nought all principality and power and virtue.
For he must reign, until he hath put all his enemies under his feet''; and
``when all things shall be subdued unto him, then the Son also himself shall
be subject unto him that put all things under him, that God may be all in
all'' (1 Cor 15:24--25, 28). Against this the Alleluia sings that his power is
\latin{potestas aeterna, quae non auferetur}, an everlasting power that shall
not be taken away, and the Communion that \latin{Sedebit Dominus Rex in
aeternum}, the Lord shall sit King for ever. Gabriel had said the same to
Mary: ``of his kingdom there shall be no end'' (Lk 1:33). If the Son delivers
the kingdom up, does he cease to reign? If he reigns ``until'' his enemies are
under his feet, does his reign stop when they are?

In St Cyril of Jerusalem's day some answered yes, and he warns his
catechumens: ``shouldest thou ever hear any say that the kingdom of Christ
shall have an end, abhor the heresy; it is another head of the dragon, lately
sprung up.'' He does not name its author; the English edition's note
identifies the teaching with Marcellus of Ancyra and his disciple Photinus.
Against it Cyril sets Gabriel's word and Daniel's vision of the Son of man,
whose power is everlasting (Dan 7:13--14), and he reads the Apostle's ``until''
as Scripture reads the word elsewhere, as setting no term: death reigned from
Adam ``until'' Moses (Rom 5:14), and the veil lies on the reader of Moses
``until this day'' (2 Cor 3:15), without either ending there. The Son's
subjection, he adds, is no new obedience, but ``a self-chosen accordance; for He
is not a servant, that He should be subjected by force, but a Son''
(\work{Catechetical Lectures} 15.27--32). St Jerome had made the same point
from Daniel: no man could be the one whose power is eternal, \latin{quae nullo
fine claudatur}, which no end encloses (\work{In Dan.}\ II, on 7:14).

\subsection{St Augustine: the faithful brought to contemplation}
\label{sec:handed-augustine}

St Augustine answers the question by asking what the kingdom is. ``The words,
When He shall have delivered up the kingdom to God, even the Father, are as
much as to say, When He shall have brought believers to the contemplation of
God, even the Father'' (\work{De Trinitate} I.8.16). The Mediator now reigns
over the just, who live by faith; at the end he brings them to sight, ``face to
face,'' and in bringing them there he is said to deliver up the kingdom. ``For
the Father will be manifested to us, and that will suffice for us,'' as Philip
asked, ``Show us the Father, and it suffices us''; and the heart is made ready
for that sight now, for ``Blessed are the pure in heart, for they shall see God''
(I.8.17). The handing over is therefore not a loss to the Son, because what he
hands over is not a possession separate from himself. ``His faithful ones,
whom He has redeemed with His blood, are called His kingdom, for whom He now
intercedes; but then, making them to abide in Himself there, where He is
equal to the Father, He will no longer pray the Father for them'' (I.10.21).

The subjection of the Son belongs to the same economy. When the kingdom is
delivered up, Christ will no longer make intercession for us ``as our Mediator
and Priest,'' and ``in so far as He is a Priest that has taken the form of a
servant for us, shall be put under Him who has put all things under Him''; in
so far as he is God, he with the Father will have put us under him (I.10.20).
Augustine reports, without adopting it, the view of those who think the human
nature in Christ will be changed into the substance of the Godhead, and
answers it by keeping the form of the servant, the Priest, as what is subject
(I.8.15; I.10.20).

On the Communion's psalm Augustine gives the session the same two times.
``The Lord inhabiteth the deluge'': ``The Lord therefore first inhabiteth the
deluge of this world in His Saints, kept safely in the Church, as in the ark.
`And the Lord shall sit a King for ever.' And afterward He will sit reigning in
them for ever'' (\work{Enarrationes in Psalmos} 28.10). The strength of the
next verse is given to his people fighting the storms of this world, ``for
peace in this world He hath not promised them''; the blessing in peace is
``affording them peace in Himself'' (28.11). The King sits in his saints now,
in the flood, and afterwards for ever; his reign in them does not stop when
they are delivered up, because he delivers them up by bringing them where he
is.

\subsection{St Hilary: not his power, but us}
\label{sec:handed-hilary}

St Hilary of Poitiers reaches the same answer by the grammar of the verse.
``He shall deliver the Kingdom to God the Father, not in the sense that He
resigns His power by the delivering, but that we, being conformed to the glory
of His body, shall form the Kingdom of God. It is not said, He shall deliver up
His Kingdom, but, He shall deliver up the Kingdom, that is, deliver up to God
us who have been made the Kingdom by the glorifying of His body'' (\work{De
Trinitate} XI.39). Those he delivers are those he has called into his kingdom,
``to whom also He has promised the blessedness of this Mystery, Blessed are the
pure in heart, for they shall see God.'' Augustine and Hilary, then, both
make the handing over the vision of God, and both cite the same beatitude for
the purity it requires.

On the subjection that follows, Hilary speaks in terms Augustine does not use.
The end of the subjection is ``that God may be all in all, that no trace of the
nature of His earthly body may remain in Him. Although before this time the two
were combined within Him, He must now become God only; not, however, by casting
off the body, but by translating it through subjection; not by losing it
through dissolutions, but by transfiguring it in glory'' (XI.40). Hilary insists
that the body is translated and not cast off, while his phrase ``God only'' goes
further than Augustine's keeping of the servant's form as what is subject.
The Preface uses neither formulation: it says only that the kingdom is handed
over and that it is eternal.

Hilary also says, expounding the Offertory's psalm, that the nations given to
the Son as inheritance are the eternal inheritance of the firstborn from the
dead (\work{Tractatus super Psalmos} 2.31); in the \work{De Trinitate} he says
that the kingdom the Son delivers to the Father is ``us who have been made the
Kingdom'' (XI.39). This reading puts the two statements together: what is
given to the Son is what the Son gives to the Father, and the Offertory's
inheritance is the kingdom that the Preface has him hand over.

\subsection{The Communion: a throne in the cleansed soul, and a throne after
the judgment}
\label{sec:handed-communion}

The Communion's psalm, which the Missal cuts down to the session and the
blessing, carries in its full verse the flood that the witnesses read in
different ways, and their readings place the session for ever at different
moments.

St Basil the Great reads it of baptism. The flood, \textit{kataklysmos}, is a
water that effaces all beneath it and cleanses all that was defiled, and it is
the grace of baptism that the psalm calls a deluge, so that the soul washed of its
sins and purged of the old man is made fit to be God's dwelling in the Spirit.
On ``And the Lord shall sit king for ever'' he says that God, settled in the soul
that shines from that deluge, makes it, as it were, a throne for himself; and
the peace of the last verse, the stability of the soul's governing part, seems
to him the most perfect of blessings, the peace the Lord gave his disciples
(\work{Homilia in Psalmum} 28.8; Jn 14:27). Cassiodorus hears the threefold
\latin{Dominus} of the two verses as the Trinity dwelling in the waters of
baptism, and the \latin{sedebit} as the Son's (\work{Expositio Psalmorum} 28).
St Thomas Aquinas, among his mystical readings, gives the world and its carnal
men brought into God's temple by conversion, after which \latin{sedebit rex in
aeternum}; strength is given for progress and the blessing in peace belongs to
the end (\work{Postilla super Psalmos} 28; Isa 32:18). Theodoret sees the world
drowned by the torrent of lawlessness rebuilt as a new creation, in which the
Lord will give impiety no place again and will reign for ever.

St Robert Bellarmine, on his second interpretation, places the session after
the judgment: when the Lord on that day has discovered the thick woods, ``then
he will `make a flood to dwell,' inundating the wicked with all manner of evils;
and thus, all resistance being broken down, the whole power of demons, bad
men, and all power in general being swept away, `the Lord shall sit King
forever.' '' The blessing in peace of the last verse is everlasting life
(\work{Explanatio}, Ps~28). His sweeping away of ``all power in general''
matches in thought the Preface's \latin{suo subiectis imperio omnibus
creaturis} and the Apostle's bringing to nought ``all principality and power,''
though Bellarmine does not cite that verse here.

Blessed Ildefonso Schuster, commenting on the antiphon as it stands in this
Mass, calls it ``a prelude to the final effect of the eucharistic Sacrifice
which ends with the divine blessing''; at the moment of Holy Communion ``Jesus is
enthroned in the faithful heart,'' and he is the Blessed through whom, by the
promise to Abraham, all nations are blessed (Gen 22:18; p.~204).

The witnesses therefore give the Communion two moments. The session has a
beginning now: in the soul cleansed by baptism and made his throne (Basil), in
the communicant (Schuster), in the saints amid the flood of this world
(Augustine). It has a completion after the judgment, Augustine's
``afterward'' and Bellarmine's sweeping away of every power. Augustine's
session runs continuously from the one into the other; Bellarmine's begins
when the judgment is done.

\subsection{How the whole Mass moves toward the Father}
\label{sec:handed-elements}

The \textbf{Preface} is the reading's text: the anointed Priest and King
subjects all creatures and hands over \latin{aeternum et universale regnum},
and the kingdom is named by goods that belong to persons, not territory:
truth and life, holiness and grace, justice, love and peace. Schuster says
that its conception of the Pontiff-King who restores the realm of creation to
the Father ``that God might eternally be all in all'' is ``borrowed from St
Paul'' (p.~204). The \textbf{Alleluia}'s \latin{potestas aeterna, quae non
auferetur} is the text Cyril of Jerusalem sets against those who said the
reign would end. The \textbf{Communion} sings the session for ever, begun in
the cleansed and fed soul and completed after the judgment, and its blessing
in peace is the peace given in himself (Augustine), the blessing of the end
(Aquinas), everlasting life (Bellarmine).

The \textbf{Epistle} gives the head who goes first and brings the body after:
\latin{primogenitus ex mortuis}, through whom the Father reconciles all things
``both as to the things that are on earth and the things that are in
heaven.'' Chrysostom says the things in heaven are reconciled because Christ
brought man up to the angels (\work{Homilies on Colossians} 3), and Theodoret
that through his sacrifice he joined the heavenly to the earthly. The
\textbf{Collect}'s \latin{omnia instaurare} is the summing up of all things in
Christ of Eph 1:10, which Chrysostom joins to Col 1:20 and Theodoret to Col
1:16; the Collect asks for its fulfilment among the nations, the work whose
completion is the kingdom handed over. The \textbf{Introit}'s \latin{in saecula
saeculorum} closes the doxology of John's greeting (Apoc 1:6), which the
antiphon sets after the angels' acclamation of the Lamb; and Schuster reads its verse as the Messianic office
conferred on Jesus, who governs ``until \dots\ he shall have conquered the
enemy'' and then lays ``the sceptre and the crown'' at his Father's feet
(pp.~201--202).

The \textbf{Gradual}'s homage of all kings and all nations is the universal
subjection that the Apostle's ``when all things shall be subdued unto him''
completes, given by some now and by all after the resurrection (Theodoret).
The \textbf{Gospel}'s \latin{regnum meum non est hinc} names a kingdom whose
origin is the Father's, and its end is the Father's too; the King bears
witness to the truth, and the vision of that truth is what Augustine and
Hilary make the handing over. The \textbf{Offertory}'s inheritance is, for Hilary,
the eternal inheritance of the firstborn from the dead; on this reading it is
given to the Son to be given by him to the Father. The \textbf{Secret} shows the
Priest still offering and still giving unity and peace now; Augustine says
that when the kingdom is delivered up he will no longer intercede for us as
Priest (\work{De Trinitate} I.10.20). And the \textbf{Postcommunion} asks for
the faithful, after the food of immortality, a share in the reign that does
not end, \latin{in caelesti sede iugiter regnare}.

\begin{fourSenses}
\item[Literal.] Daniel sees one like a son of man given an everlasting
kingdom; Paul says that the Son, firstborn from the dead, reconciles all
things, and elsewhere that he will deliver up the kingdom to the Father; the
psalm says that the Lord makes the flood to dwell and sits King for ever,
which the witnesses read of God's judgments from the beginning (Aquinas's
plain sense) or, mystically, of the world, of the Church in the ark, and of
baptism.
\item[Allegorical.] Christ the head brings his body, the Church, to the
Father, and the Church is the kingdom handed over (Augustine, Hilary). The
flood is baptism, out of which the soul becomes his throne (Basil).
\item[Moral.] The present life is the purification of heart that the vision
requires, ``Blessed are the pure in heart,'' cited by both Augustine and
Hilary; it is lived in the peace he gives in himself rather than the world's,
and the soul cleansed in baptism and fed at the altar is made his throne
(Basil; Schuster).
\item[Anagogical.] God all in all, and the vision face to face: an endless
reign of the redeemed with Christ, who is not diminished by giving them to the
Father, after the judgment has swept away every power that resisted him
(Bellarmine).
\end{fourSenses}
